% commentaire.tex — f162v, vv. 1–42 (conversion initiale des travaux antérieurs)
\comm{1–10}{ὦ Ζεῦ ... οὐχ ἅ γ᾽ ἂν ἀκούειν δέῃ.}{Nous sommes au tout début de la pièce, il fait encore nuit, un personnage entre d'un pas pressé ; il est immédiatement reconnaissable à son masque bien connu des spectateurs : c'est Euripide\nc{À quoi pouvait ressembler ce masque d'Euripide ? On peut supposer qu'il présentait une barbe très touffue (cf. le vers 190) ainsi que de multiples taches sur son visage : Ἐλέγετο δὲ καὶ βαθὺν πώγωνα θρέψαι καὶ ἐπὶ τῆς ὄψεως φακοὺς ἐσχηκέναι. (Tiré d'une des \italique{Vies} d'Euripide, rassemblées par Richard Kannicht pour son édition des fragments d'Euripide, cinquième volume des \italique{Tragicorum Graecorum Fragmenta} (2004). Ce détail se trouve à la page 47.)}. Il est suivi d'un vieillard qui traîne la patte et crie sa détresse -- \italique{topos} comique. (On apprendra plus loin qu'il est un parent d'Euripide ; la tradition l'identifie à Mnésiloque, le père de la première femme d'Euripide, Chœrile\nc{\italique{Confer} l'introduction sur la vie d'Euripide et de Mnésiloque et le vers 291.}.) Réplique fracassante d'Euripide, et bien selon ses façons -- les fameuses γνῶμαι d'Euripide fort goûtées des spectateurs : ἀλλ᾽ οὐκ ἀκούειν δεῖ σε πάνθ᾽ ὅσ᾽ αὐτίκα ὄψει παρεστώς, « il n'est pas besoin que tu entendes tout ce que tu vas bientôt voir ».
Hors d'haleine, Mnésiloque n'est pas en état de suivre Euripide dans ses aphorismes : il s'est arrêté à la première partie : οὐ δεῖ μ᾽ ἀκούειν, « il ne faut pas que j'entende... » Euripide reprend : οὐχ ἅ γ᾽ ἂν μέλλῃς ὁρᾶν, « pas certes ce que tu es sur le point de voir ».\\
C'est là que nous arrivons à l'étrange réponse de Mnésiloque au vers 7b : οὐδ᾽ ἆρ᾽ ὁρᾶν δεῖ μ᾽, « je ne dois pas non plus voir ? ». Ne comprenant pas pourquoi  Mnésiloque passe d'une défense d'entendre à une défense de voir, l'on est incité à mieux scruter les vers précédents. 
Le scholiaste nous met sur la piste : les constructions tortueuses d'Euripide sont sujettes à l'amphibologie\nc{ὁ μὲν τραγικώτερον καὶ ὑψηλότερον φράζει, ὁ δὲ ταπεινότερον ἀκούει, « l'un s'exprime dans un style plus tragique et élevé, quand l'autre comprend de façon plus basse ».}.
Mnésiloque commet une erreur (volontaire) d'interprétation sur la dernière réplique d'Euripide : dans `οὐχ ἅ γ᾽ ἂν μέλλῃς ὁρᾶν', il feint d'entendre `ἄγαν' à la place de `ἅ γ᾽ ἂν'. Aristophane ayant mis à dessein ces particules dans un ordre inhabituel\nc{Le véritable Euripide s'exprimait ainsi : (\italique{Oreste}, vers 81) : Ἑλένη, τί σοι λέγοιμ᾽ \textbf{ἂν ἅ γε} παροῦσ᾽ ὁρᾷς;}.
Il est remarquable de noter que, des onze pièces conservées de notre poète, le cas ne se rencontre qu'une seule autre fois... et c'est encore un pastiche d'un Euripide postiche (celui des \italique{Grenouilles}, cette fois, vers 960) :\citComm{ἐξ \textbf{ὧν γ᾽ ἂν} ἐξηλεγχόμην.}Mnésiloque croit comprendre que le voir est lui aussi soumis au μηδὲν ἄγαν grec : `οὐκ ἄγαν μέλλῃς ὁρᾶν', « tu ne dois pas trop voir »... d'où sa réaction du vers 7b.\\}
\comm{1–4}{ὦ Ζεῦ ... πρὶν τὸν σπλῆνα κομιδῇ μ᾽ ἐκβαλεῖν.}{Ce quatrain est analogue à celui des vers 71-74. À deux reprises, Mnésiloque adresse une prière plaintive à Zeus avant de questionner Euripide. Ressemblance saisissante aussi avec les premiers vers des \italique{Nuées} (vv. 1-5) :\citComm{Στρεψιάδης --- ἰοὺ ἰού·\\
ὦ Ζεῦ βασιλεῦ τὸ χρῆμα τῶν νυκτῶν ὅσον·\\
ἀπέραντον. οὐδέποθ᾽ ἡμέρα γενήσεται;\\
καὶ μὴν πάλαι γ᾽ ἀλεκτρυόνος ἤκουσ᾽ ἐγώ·\\
οἱ δ᾽ οἰκέται ῥέγκουσιν· ἀλλ᾽ οὐκ ἂν πρὸ τοῦ.\\
Strepsiade --- Iou ! iou !\\
Ah ! Zeus roi, quel fardeau que ces nuits !\\
C'est interminable... Il ne fera donc jamais jour ?\\
Pourtant voilà longtemps que j'ai entendu le coq,\\
et mes serviteurs ronflent. Ah ! que non pas autrefois !}
Cette analogie permet de saisir le sens imagé de l'hirondelle (χελιδών) de notre passage : elle est mise à la place du coq des \italique{Nuées} pour ce qu'elle a de semblable : en chantant dès avant l'aube, elle annonce le jour bientôt renaissant\nc{Et non pas le printemps -- qui n'est pas de saison ici --, comme on peut le lire drôlement dans nombre de commentaires.}. Que l'hirondelle était réputée chez les Grecs chanter\nc{Chanter ?.. ou plutôt babiller : l'hirondelle était proverbialement brocardée pour son caquet, cf. notamment Taillardat §§520-522 et les notes d'Alphonse Willems à propos du vers 682 des \italique{Grenouilles}.} au matin, cela est appuyé par l'adjectif ὀρθρογόη, \italique{aux lamentations matinales}, qu'Hésiode (\italique{Opera et Dies}, v. 568) applique à l'hirondelle\nc{Voir aussi Hésychius (ο 1189) : ὀρθ<ρ>ογόη· ἡ χελιδών.} :\citComm{ὀρθρογόη Πανδιονὶς ὦρτο χελιδὼν\\
La fille matinale et gémissante de Pandion, l'hirondelle, réveille...}Ce strident réveille-matin avait donné lieu à une expression proverbiale dont Aristophane fait mention ailleurs\nc{Fragment 617 Kassel-Austin. Diversement recensé par les parémiographes (souvent sous la forme πυθοῦ Χελιδόνος) : Hésychius π 4311 ; Photius π 1519 ; Suidas π 3136 ; \Apostol XV, 14 ; \&c.}:\citComm{πυθοῦ χελιδὼν πηνίκ᾽ ἄττα φαίνεται\\
Renseigne-toi à quelle heure l'hirondelle apparaît.}Ici, ce n'est pas une hirondelle mais Euripide qui a jeté hors du lit Mnésiloque bien trop tôt.. Ce sera donc lui qui sera interrogé : \textbf{παρὰ σοῦ} πυθέσθαι (v. 3).}
\comm{2}{ἀλοῶν.}{Ἀλοᾶν est un terme agricole signifiant \italique{battre le grain}, selon une technique bien définie : sur l'aire à battre (ἡ ἅλως), on fait tourner des bœufs, attachés à un piquet, qui piétinent et qui broient la céréale jonchant le sol. De ce fait, ἁλοᾶν est devenu un synonyme expressif de \grec{τύπτειν} avec le sens de \italique{broyer de coups}. (Note de Taillardat\nc{Je me référerai souvent – et avec force citations ! – au magistral ouvrage que Jean Taillardat a consacré à la langue bigarrée et haute en couleurs d'Aristophane : \italique{Les images d'Aristophane} édité aux Belles Lettres (seconde édition de 1965). Je cite ici le §138 ; voir aussi les références du §606 : la métaphore était usuelle et pas de l'invention de notre poète : \italique{Grenouilles} 149, Sophocle fr. 19 \Nauck [$=$ fr. 20 Radt (dub.)] ; Hérondas 2, 34 ; Babrius, \italique{Le lion amoureux}, 98.15 ; καταλοᾶν Eschine 2, 140 ; Xénoph. \italique{Cyrop.} 7, 1, 31 ; Lucien \italique{Icarom.} 15 ; συναλοᾶν Hippon. 14AD\expo{3}.4 [$=$ fr. 92.4 Masson $=$ fr 95.4 Degani] ; Théocr. 22.128.})\\
Ici, ἁλοᾶν a le sens moral de \italique{tourmenter}. Conformité du grec et du français -- ou du latin, plutôt : l'on possède la même image avec notre \italique{tribulation}, le latin \latin{tribulare} signifie « battre le blé, presser avec la herse (\latin{tribulum}) ».}
\comm{3}{πρὶν τὸν σπλῆνα κομιδῇ μ᾽ ἐκβαλεῖν.}{Les dérangements (palpitations) dus à la course – d'autant plus aggravés par l'âge avancé de Mnésiloque – étaient imputés par les Anciens à des troubles de la rate (σπλήν, \anglais{spleen} en anglais)\nc{Voir Pline l'Ancien, \italique{Histoire naturelle} 11,205 ; Plaute, \italique{Mercator}, v. 124.}. En français, celui qui court ventre à terre court aussi comme un « dé-raté ».}
\comm{13–8}{}{Aristophane nous peint Euripide comme un de ces faiseurs de cosmogonies (tel le Socrate des \italique{Nuées} avec son Δῖνος). Le véritable Euripide avait effectivement un penchant pour ce genre de nuées\nc{Voir par exemple le fragment 5 de \italique{Mélanippe la sage} dans l'édition de Jouan-van Looy aux Belles Lettres, tome VIII, \nieme{2} partie ($=$ fragment 484 de l'édition Kannicht).}. En l'occurrence, le pastiche de cosmogonie qu'Aristophane met dans la bouche d'Euripide semble inspiré d'Empédocle\nc{Cf. notamment ce passage d'Eusèbe de Césarée (\italique{Préparation évangélique}, XV, 37 ; fragment A49 Diels-Kranz) : Ἐμπεδοκλῆς τὸν μὲν αἰθέρα πρῶτον διακριθῆναι et la création de l'œil par Aphrodite (B84.9 D.-K.) reconstitué par Blass (1883) : χοάνῃσι δίαντα τετρήατο θεσπεσίῃσιν.}.}
\comm{16–8}{πρῶτα ἐμηχανήσατο ... ἀκοὴν δὲ χοάνης ὦτα διετετρήνατο.}{Euripide recourt à une image tirée de la fonderie -- certainement plus parlante aux spectateurs de l'époque qu'à nous : l'oreille (ἀκοή) est comparée aux ὦτα (littéralement les \italique{oreilles} et, par analogie, les \italique{anses}, \italique{poignées} ou, comme ici, les \italique{godets} ou \italique{entonnoirs} d'un moule, χοάνη). Pour le texte, voir le commentaire critique.}
\comm{19}{διὰ τὴν χοάνην οὖν μήτ᾽ ἀκούω μήθ᾽ ὁρῶ;}{Mnésiloque ne goûte guère aux subtilités des poètereaux et tord à plaisir leurs propos. Ses saillies fort désobligeantes se repèrent aisément puisqu'elles sont toujours superbement ignorées. C'en est ici le premier exemple ; Euripide le-faiseur-de-cosmogonies euripidise en lui racontant des histoires de fonderie et de moule... conclusion de Mnésiloque pour expliquer son accablement : « C'est donc à cause de ton moule que je suis moulu ? », διὰ τὴν χοάνην οὖν μήτ᾽ ἀκούω μήθ᾽ ὁρῶ\nc{Ou bien : « Je ne vois plus Claire ; j'ai perdu Louis. » (traduction plus libre.)}. C'est bien le sens qu'il faut donner au tour `μήτ᾽ ἀκούω μήθ᾽ ὁρῶ' et qu'on peut rapprocher du latin \citComm{\latin{Nec aures me credo habere nec tango},\\
Je crois n'avoir plus ni oreilles ni main,} qui signifie justement qu'on ne se sent plus la force de faire quoi que ce soit\nc{Je renvoie le lecteur à l'admirable \italique{Dictionnaire des sentences latines et grecques} de Renzo Tosi (traduit de l'italien par Rebecca Lenoir ; proverbe n° 1971, cité par saint Jérôme notamment). Érasme avait recensé le proverbe dans ses \italique{Adages}, n° 2500. (Je cite les \italique{Adages} ainsi que leur traduction française selon l'édition de Jean-Christophe Saladin aux Belles Lettres (2019).) Toujours dans le §1971, Tosi cite aussi la sentence médiévale \latin{Nec audio nec video}, qui se superpose parfaitement au grec μήτ᾽ ἀκούω μήθ᾽ ὁρῶ.}.\\
La langue d'Aristophane, comme nous l'allons voir en chemin, est truffée d'échos aux proverbes grecs ; on ne peut que regretter le manque d'intérêt criant des commentateurs d'Aristophane pour ce domaine. Je m'efforcerai au contraire, et selon mes moyens, de souligner chacun des emprunts d'Aristophane à cette langue proverbiale. Et ils sont nombreux ! Le cas présent est malheureusement moins probant puisque je n'ai pas pu trouver d'autres occurrences dans la littérature grecque. (Surtout : le jeu de mots exact en grec, qui tourne autour de χοάνη, nous échappe.)}
\comm{21}{οἷόν τέ που 'στιν αἱ σοφαὶ ξυνουσίαι.}{Le scholiaste note à ce vers une allusion à un mot célèbre : σοφοὶ τύραννοι τῶν σοφῶν ξυνουσίᾳ, qu'Aristophane dans ses \italique{Héros}\nc{Fragment 323 de l'édition Kassel-Austin.} mais aussi Platon\nc{\italique{Théagès}, 125b et la \italique{République}, VIII, 568b. Platon devait avoir un point de vue sur Euripide assez proche de celui d'Aristophane : il chasse joyeusement de sa République idéale Euripide, cet homme qui occupe une place prestigieuse parmi les semi-habiles de sa cité.} et Antisthène\nc{Fragment 59 de l'édition Decleva-Caizzi} attribuent à Euripide. Cela n'empêche pas moins notre scholiaste, bien sûr de lui, d'accorder la paternité de cette formule à Sophocle\nc{(Dans son Αἴας Λοκρός.) Fr. 14 de l'édition de Stefan Radt, \italique{Tragicorum Graecorum Fragmenta}, vol. 4, édition corrigée et augmentée de 1999. Il figure aussi parmi les fragments incertains d'Euripide (n° 888b) dans l'édition Kannicht des \italique{Tragicorum Graecorum Fragmenta}, vol 5. Le ton de la scholie est d'une telle assurance doublée d'une forte érudition qu'on serait tenté de l'en croire. Preuve que la scholie est de Didyme, selon le critère de Pierre Boudreaux ? (\italique{Le texte d'Aristophane et ses commentateurs} (1920). Boudreaux, brillant helléniste fauché bien jeune par la Grande guerre, réussit le tour de force d'écrire un bouquin -- publié de façon posthume -- passionnant quoique traitant de sujets à priori fort arides.)
Après tout, Euripide est coutumier du fait : Aulu-Gelle (\italique{Nuits attiques}, XIII, 18) nous rapporte entre autres exemples que le vers 193 des \italique{Bacchantes}, γέρων γέροντα παιδαγωγήσω σ᾽ ἐγώ, se retrouve aussi à l'identique chez Sophocle (fr. 695 Radt) ; de même, Aristote (Περὶ ποιητικῆς, XIII, 1458b22) cite une reprise d'Eschyle par Euripide.
À rebours de leur querelle dans les \italique{Grenouilles}, Aristote cite un vers iambique d'Eschyle jugé plat en raison du recours à un mot courant tandis qu'Euripide le reprit en l'agrémentant d'un mot emprunté et rare : οἷον τὸ αὐτὸ ποιήσαντος ἰαμβεῖον Αἰσχύλου καὶ Εὐριπίδου, ἓν δὲ μόνον ὄνομα μεταθέντος, ἀντὶ κυρίου εἰωθότος γλῶτταν, τὸ μὲν φαίνεται καλὸν τὸ δ᾽ εὐτελές. Αἰσχύλος μὲν γὰρ ἐν τῷ Φιλοκτήτῃ ἐποίησε·\textbf{φαγέδαιναν ἥ μου σάρκας ἐσθίει ποδός,} ὁ δὲ ἀντὶ τοῦ `ἐσθίει' τὸ `θοινᾶται' μετέθηκεν.}. Toujours est-il qu'Euripide n'en passait certainement pas moins aux yeux d'Aristophane et de ses spectateurs pour le père putatif de la citation. Sur le sens de γέ που, on lira avec intérêt la seconde note de Willems reproduite ci-après en annexe.}
\comm{22–4}{πῶς ἂν οὖν ... προσμάθοιμι χωλὸς εἶναι τὼ σκέλει;}{Aristophane joue sur un détournement de proverbe équivalent à notre «~Qui fréquente le boiteux apprend à boiter~» : Παρὰ χωλὸν οἰκῶν κἂν ἐπισκάζειν μαθήσῃ\nc{C'est le proverbe n° 1674 du \italique{Dictionnaire des sentences latines \& grecques}. Recensé par Macaire au \nieme{XIV} siècle dans son anthologie Ῥοδωνιά (VI, 90 de l'édition de Leutsch \& Schneidewin du \italique{Corpus Paroemiographorum Graecorum}, vol. II) ainsi que sous d'autres formes proches par Michel Apostolius, Plutarque et un scholiaste de Pindare.}.
Le boiteux -- ou plutôt le \italique{faiseur de boiteux}, χωλοποιός\nc{\italique{Les Grenouilles}, vers 846.} -- c'est bien sûr Euripide et sa fâcheuse prédilection pour ce genre de personnages\nc{Ou autres aussi peu reluisants. Cf. les \italique{Acharniens} v. 411 sq. : Télèphe, Bellérophon, Philoctète le mendiant ; voir aussi la \italique{Paix} v. 147 ; les \italique{Oiseaux} v. 1293 et le fragment 53 de l'édition Kassel-Austin des \italique{Poetae Comici Graeci}, vol. III 2 (1984).}.
La pointe de la plaisanterie de Mnésiloque est rejetée à la toute fin ; quand il précise vouloir être boiteux τὼ σκέλει, \italique{des deux jambes}. 
Le scholiaste a raison : \citComm{χωλὸς εἶναι· ἵνα μὴ περιπατήσῃ μετ᾽ αὐτοῦ.}Si Mnésiloque venait à boiter des deux jambes, il ne serait plus forcé de suivre Euripide. Le grec aime à préciser que l'on est rompu \italique{des deux jambes} pour souligner une fatigue intense\nc{Voir par exemple la \italique{Paix} et la douloureuse ascension de l'orchestre à la scène de Trygée : ἔγωγέ τοι πεπόνηκα κομιδῇ τὼ σκέλει (v. 820) et encore : ἤλγουν τὼ σκέλει μακρὰν ὁδὸν διεληλυθώς (v. 825). 
Nous ne sommes pas dans la chanson de Nadaud où souhaiter claudiquer des deux côtés, c'est souhaiter ne plus claudiquer :\citComm{Si je marche plus droit qu'un autre\\
C'est que je boite des deux pieds.}}.
Mais Euripide ne s'en laisse pas compter et lui ordonne de lui emboîter le pas (βάδιζε δευρὶ).\\
Pour le texte, voir le commentaire critique.}
\comm{27–8}{σίγα νυν --- σιωπῶ τὸ θύριον --- ἄκου᾽ --- ἀκούω καὶ σιωπῶ τὸ θύριον.}{Le trait d'esprit de Mnésiloque nous échappe.}
\comm{29}{Ἀγάθων.}{Voici enfin la réponse à la question de Mnésiloque (v. 4) sur la destination de leur équipée : la demeure d'Agathon, le fameux poète tragique\nc{On n'a conservé d'Agathon que des fragments et quelques témoignages. Ils ont été regroupés par Snell \& Kannicht au premier volume des \italique{Tragicorum Graecorum Fragmenta} (1986).} – sauf pour Mnésiloque –, rendu célèbre par sa victoire au concours tragique des Lénéennes de 416\nc{Athénée, \italique{Deipnosophistes}, V, 216f-217a.}. Platon immortalisa dans son \italique{Banquet} la réception qu'il fit chez lui le lendemain de cette victoire. Le rôle d'Agathon dans les \italique{Thesmophories} est semblable à celui d'Euripide dans les \italique{Acharniens} (v. 410-489) : il fournit le héros en matériel de travestissement. Euripide confirmera cette parenté aux vers 173-174\nc{Aristophane n'est d'ailleurs pas le seul à rapprocher Euripide \& Agathon : ainsi fait Aristote dans sa \italique{Poétique} (1456a25-32) : Agathon introduisit des interludes (ἐμβόλιμα) chantés par le chœur sans rapport avec l'action, suivant la voie tracée par Euripide.}. Après la mort d'Agathon, Aristophane lui rendit un hommage posthume dans les \italique{Grenouilles} aux vers 83-85.}
\comm{37}{μυρρίνας}{désigne les \italique{rameaux de myrte} que l'on brûlait lors des sacrifices (cf. les \italique{Guêpes}, v. 860 ; les \italique{Oiseaux} au v. 43).}
\comm{39–62}{Passage en dimètres anapestiques.}{(Ceux aux vers 42 et 62 sont catalectiques.) Ce mètre est d'usage courant chez Aristophane pour les prières (cf. les \italique{Guêpes} v. 323-333 et la \italique{Paix} v. 947-1015). Le serviteur d'Agathon est tout empreint de la préciosité du maître, comme, du reste, le serviteur de l'Euripide des \italique{Acharniens} (v. 396-400) l'était.}
\comm{39–40}{εὔφημος πᾶς ἔστω λέως / στόμα συγκλῄσας·}{Ce genre de formule est classique au début d'une prière ou d'un sacrifice. On peut néanmoins constater ici un peu plus d'emphase avec le tour tragique\nc{C'est ce que souligne le scholiaste : παρὰ τοῖς τραγικοῖς οὕτω λέγεται, ``<εὔφημος> πᾶς ἔστω λαός''.} et l'utilisation de l'adjectif εὔφημος plutôt que le classique verbe εὐφημεῖν\nc{\italique{Acharniens} v. 241 ; \italique{Cavaliers} v. 1316 εὐφημεῖν χρὴ καὶ στόμα κλείειν ; \italique{Nuées} v. 263 ; \italique{Paix} v. 434 et 1316 ; \italique{Grenouilles} v. 354.} ou substantif εὐφημία\nc{\italique{Guêpes} v. 868 ; \italique{Oiseaux} v. 959.}.}
\comm{40}{ἐπιδημεῖ.}{\italique{Séjourner}, tel me paraît bien le sens de ἐπιδημεῖ, \latin{peregre advenit et adest}. (Xénophon, \italique{Mémor.}, I, 2, 61.) Cf. ἐπιχωριάζει, \latin{venit ibique commoratur}. Μελοποιῶν a le sens du futur (voir note [de Willems] sur les \italique{Grenouilles} au v. 148, t. III p. 109-110). (Note de Willems.)}
\comm{41–2}{θίασος Μουσῶν ἔνδον μελάθρων / τῶν δεσποσύνων μελοποιῶν.}{Pour bien goûter ces vers, il faut savoir ce qu'est un \italique{thiase} : une bruyante troupe de gens célébrant des rites (bachiques la plupart du temps) et parcourant les rues. Cf. les \italique{Bacchantes} d'Euripide où des femmes munies de tambourins interprètent des danses scandées de cris rituels (évohé) avec une fougue poussée jusqu'à la frénésie. Par extension, ce mot très prisé des tragiques désigne une foule extrêmement bruyante ; Euripide l'emploie à propos d'une bande de centaures (\italique{Iphigénie à Aulis}, v. 1059) ; une troupe endiablée de jeunes danseuses (\italique{Iphigénie en Tauride}, v. 1146) ou encore une meute de guerriers (\italique{Phaéton}, v. 796). 
Pour décrire cette multitude de Muses, au lieu de recourir à la banale image de l'\italique{essaim}\nc{Le Socrate des \italique{Nuées} (vers 297) parle d'un « essaim de déesses », μέγα γάρ τι θεῶν κινεῖται σμῆνος. Voir Jean Taillardat, \italique{Les images d'Aristophane}, §662.}, à dessein Aristophane préfère évoquer un thiase : comme le remarque Sommerstein, l'inversion est totale : au lieu d'un cortège célébrant une divinité, c'est un thiase de Muses qui vient honorer Agathon.}
